%======================================================================
%  2. Project Objectives: questions (and sub-questions) answered
%     by the visualizations.
%======================================================================
\section{Project Objectives}
\label{sec:objectives}

The dashboard is built to answer the following questions.

\begin{enumerate}
  \item \textbf{How did house prices and income evolve over time?}
    \begin{enumerate}
      \item What is the trend of house prices in the different Danish provinces
            (\emph{landsdele})?
      \item How did affordability evolve for couples compared to singles,
            both for Denmark as a whole and per province?
    \end{enumerate}
  \item \textbf{Is it better to live in areas with a zip code above 5000 or not
        (DK1 vs.\ DK2)?}
  \item \textbf{Is it worth living in Copenhagen?} We expect average income to be
        higher there, but house prices to be higher as well.
\end{enumerate}

\noindent TODO --- Add, for each question, what you would like to learn and
accomplish, and which visualization(s) in Section~\ref{sec:visualization}
address it.

Intro text
We expect that provinces with higher housing prices also have higher incomes, but that housing prices have risen faster than incomes since 1995, so housing has become less affordable over time. Instead of testing this directly, we use it as a starting point and explore it through the following questions. Throughout the report, we measure affordability as the price-to-income ratio: the average house price divided by the average yearly income. A higher ratio means housing is less affordable.

Questions

How have house prices and income evolved in Denmark since 1995?
1.1 What is the trend of house prices in the different provinces?
1.2 How has affordability changed for couples compared to singles, in Denmark as a whole and in the different provinces?
Is housing more affordable east or west of the Great Belt (zip codes below or above 5000)?
Is it worth living in Copenhagen? Is the affordability gap between Copenhagen and the other provinces growing over time?
